% =====================================================================
%  THÈME 2 — Conversions en quantité de matière
%  Niveau DIFFICILE — Exercices
% =====================================================================
\documentclass[11pt,a4paper]{article}
\usepackage[utf8]{inputenc}
\usepackage[T1]{fontenc}
\usepackage{lmodern}
\usepackage[french]{babel}
\usepackage{amsmath,amssymb,mathtools}
\usepackage{icomma}
\usepackage[version=4]{mhchem}
\usepackage{siunitx}
\sisetup{output-decimal-marker={,},per-mode=symbol}
\DeclareSIUnit\Molar{\mole\per\litre}
\usepackage{xcolor}
\usepackage{tikz}
\usepackage[most]{tcolorbox}
\usepackage{enumitem}
\usepackage{array}
\usepackage{fancyhdr}
\usepackage[a4paper,margin=2.1cm,headheight=26pt,headsep=8pt]{geometry}

\definecolor{primary}{RGB}{142,93,168}
\definecolor{primarydark}{RGB}{82,45,108}
\definecolor{difcol}{RGB}{176,40,60}
\newcommand{\gmol}{\si{\gram\per\mole}}
\newcommand{\Vm}{V_{\mathrm m}}

\pagestyle{fancy}\fancyhf{}
\fancyhead[L]{\small\textcolor{primarydark}{\textbf{Fiche 5} — Thème 2 : Conversions}}
\fancyhead[R]{\small\textcolor{difcol}{\textsc{Difficile}}}
\fancyfoot[C]{\small\thepage}
\renewcommand{\headrulewidth}{0.8pt}
\renewcommand{\headrule}{\hbox to\headwidth{\color{primary}\leaders\hrule height \headrulewidth\hfill}}

\newtcolorbox[auto counter]{exo}[1][]{enhanced,breakable,boxrule=1pt,arc=3pt,
  left=3mm,right=3mm,top=2.4mm,bottom=2.4mm,colback=difcol!5,colframe=difcol,
  fonttitle=\bfseries,coltitle=white,
  title={Exercice~\thetcbcounter\ \ifx\relax#1\relax\relax\else\textmd{--- \itshape #1}\fi}}

\setlength{\parindent}{0pt}
\setlength{\parskip}{3pt}

\begin{document}

\begin{center}
{\LARGE\bfseries\color{primarydark}Thème 2 — Conversions}\\[1pt]
{\large\color{difcol}\textbf{Niveau Difficile}}\\[2pt]
{\itshape Objectif : identifier une espèce par sa masse molaire, mener des chaînes à plusieurs conversions.}\\
{\color{primary}\rule{0.6\linewidth}{1pt}}
\end{center}

\begin{exo}[identifier un métal alcalin par la masse molaire]
On dispose de \SI{0,20}{\mole} d'un phosphate d'un métal alcalin, de formule \ce{X3PO4}, dont la
masse est \SI{42,4}{\gram}. \emph{Données (\gmol) :} \ce{P}=31 ; \ce{O}=16 ; et pour les alcalins
\ce{Na}=23 ; \ce{K}=39 ; \ce{Rb}=85{,}5 ; \ce{Cs}=133.
\begin{enumerate}[label=\textbf{\alph*)},leftmargin=2em,topsep=1pt,itemsep=1.5pt]
  \item Quelle est la masse molaire $M$ du phosphate \ce{X3PO4} ?
  \item Exprime $M(\ce{X3PO4})$ en fonction de la masse molaire $M_{\ce{X}}$ du métal alcalin.
  \item Résous pour obtenir $M_{\ce{X}}$.
  \item Identifie le métal \ce{X}.
\end{enumerate}
\end{exo}

\begin{exo}[un sulfate métallique hydraté (mélantérite)]
La mélantérite est un sulfate métallique heptahydraté de formule \ce{MSO4.7H2O}. On mesure que
\SI{0,20}{\mole} de ce solide pèsent \SI{55,6}{\gram}. \emph{Données (\gmol) :} \ce{S}=32 ;
\ce{O}=16 ; \ce{H}=1 ; et pour \ce{M} : \ce{Co}=59 ; \ce{Fe}=56 ; \ce{Cu}=63{,}5 ; \ce{Zn}=65{,}4.
\begin{enumerate}[label=\textbf{\alph*)},leftmargin=2em,topsep=1pt,itemsep=1.5pt]
  \item Calcule la masse molaire du solide hydraté.
  \item Sachant que l'eau de cristallisation compte pour $7\times\SI{18}{\gram\per\mole}$, exprime la
        masse molaire en fonction de $M_{\ce{M}}$ et résous.
  \item Identifie le métal \ce{M} (et nomme le composé).
\end{enumerate}
\end{exo}

\begin{exo}[du cube solide aux ions en solution]
On considère un cube d'hydroxyde de sodium \ce{NaOH} d'arête $a=\SI{2}{\centi\meter}$ et de masse
volumique $\rho=\SI{2,0}{\kilo\gram\per\deci\meter\cubed}$. \emph{Donnée :}
$M(\ce{NaOH})=\SI{40}{\gram\per\mole}$.
\begin{enumerate}[label=\textbf{\alph*)},leftmargin=2em,topsep=1pt,itemsep=1.5pt]
  \item Calcule le volume du cube (en \si{\centi\meter\cubed}) et convertis la masse volumique en
        \si{\gram\per\centi\meter\cubed}.
  \item En déduire la masse de \ce{NaOH}.
  \item Combien de moles de \ce{NaOH} cela représente-t-il ?
  \item En déduire la quantité de matière d'\textbf{ions sodium} \ce{Na+} (une \ce{NaOH} libère un \ce{Na+}).
\end{enumerate}
\end{exo}

\begin{exo}[deux ions dans une eau minérale]
Une étiquette d'eau minérale indique \SI{90,0}{\milli\gram\per\litre} d'ions calcium \ce{Ca^2+} et
\SI{122}{\milli\gram\per\litre} d'ions hydrogénocarbonate \ce{HCO3-}. \emph{Données (\gmol) :}
\ce{Ca}=40 ; \ce{H}=1 ; \ce{C}=12 ; \ce{O}=16.
\begin{enumerate}[label=\textbf{\alph*)},leftmargin=2em,topsep=1pt,itemsep=1.5pt]
  \item Calcule la concentration molaire en ions \ce{Ca^2+} (en \si{\Molar}).
  \item Calcule la masse molaire de l'ion \ce{HCO3-}, puis sa concentration molaire.
\end{enumerate}
\end{exo}

\end{document}
